% ---------------------------------------------------------------------------
% Experiment History -- new subsection to insert IMMEDIATELY AFTER the scope
% paragraph of \section{Experiment History} and BEFORE
% \subsection{Simulation Chain and Validation}.
%
% Purpose: state the objective once, in one place, so that every later
% subsection can name its delta in one clause instead of restating a formula.
% The table doubles as the reproducibility artefact for the findings paper.
%
% ENDS AT THE TABLE, deliberately. The three paragraphs that used to follow
% were cut 2026-09-10: their two findings belong at their point of use, not
% here, and their third paragraph duplicated \ref{sec:system:objective}.
%   - log barrier -log(1-p_hat+eps) is ~constant across the batch when the
%     detector flags everything, so normalised advantages are noise
%     -> goes to subsection 6 (untrainability), where it is the mechanism.
%   - gamma = 0.02 never binds, so the two-agent result is an
%     unconstrained-power result
%     -> goes to subsection 4 (two agents), where the number is read.
%   - "power is a constraint, not a preference" -> already stated in
%     System Model, \ref{sec:system:objective}. Do not restate it here.
%
% All constants re-read from the code on 2026-09-10:
%   simulation01/jammer_env.py:10,72     BETA=3,    reward = ber - BETA*flagged - 0.05
%   simulation02/jammer_env.py:10,74     BETA=0.5,  same shape, kurtosis detector
%   simulation03/jammer_env.py:9,10,80   BETA=2.0,  GAMMA=0.02, hinge on kurtosis excess
%   simulation03c/jammer_env.py:9,10,80  identical to sim03
%   simulation03b/train.py:31,32,234,249 LAMBDA=2.0, GAMMA=0.02, loss form
%   simulation04/train.py:31,32,288,290  LAMBDA=2.0, GAMMA=0.02, power term on the SUM
%   simulation06/train_jammer.py:50,51,321-323  BETA=0.3 (warmup), GAMMA=0.1, log barrier
%   simulation07/*.py:68,69,437-439      BETA=0.3 (warmup), GAMMA=0.05
% ---------------------------------------------------------------------------

\subsection{The Attacker's Objective}
Every attacker below optimises the same three quantities -- how many bits it breaks, how often it is
flagged, and how much power it spends -- and the experiments differ in how those three are combined.
Stating the combination once makes the rest of this appendix legible, because several of the results
that follow are properties of the objective rather than of the attacker.

In the reinforcement-learning experiments the per-frame reward is
\begin{equation}
R \;=\; \mathrm{BER} \;-\; \beta \, D(\mathbf{y}) \;-\; \gamma \, P \;-\; c ,
\label{eq:reward}
\end{equation}
with $P$ the mean jammer power over the frame, $c = 0.05$ a small constant so that transmitting
nothing is not free, and $D$ a detection term whose form changed three times.
Where the defender is the kurtosis test, $D$ is written in terms of $\kappa(\mathbf{y})$, the Fisher
excess kurtosis of the received frame averaged over its in-phase and quadrature components; it is
close to $-2$ for a clean QPSK frame and to $0$ for Gaussian noise, and the detector flags a frame
when it exceeds $\tau = -1.0$.
In the directly optimised experiments the same three quantities appear as a loss,
\begin{equation}
\mathcal{L} \;=\; \mathcal{L}_{\mathrm{BER}} \;+\; \lambda \, D(\mathbf{y}) \;+\; \gamma \, P ,
\label{eq:loss}
\end{equation}
where $\mathcal{L}_{\mathrm{BER}}$ is a differentiable surrogate for the bit error rate: the binary
cross-entropy between the demapper's log-likelihood ratios and the \emph{complement} of the
transmitted bits, so that driving the loss down drives the error rate up.
Table~\ref{tab:objectives} gives the instantiation used at each step.

\begin{table*}[!t]
\renewcommand{\arraystretch}{1.3}
\caption{The attacker's objective at each step. $D$ is the detection term of \eqref{eq:reward} and \eqref{eq:loss}; $\kappa(\mathbf{y})$ is the Fisher excess kurtosis of the received frame, averaged over its in-phase and quadrature components; $\tau=-1.0$ is the detector's threshold on it; and $\hat{p}$ is the learned detector's output probability.}
\label{tab:objectives}
\centering
\begin{tabular}{l l l c c l}
\hline\hline
\textbf{Step} & \textbf{Detector} & \textbf{Detection term $D$} & \textbf{$\beta$ / $\lambda$} & \textbf{$\gamma$} & \textbf{Optimised by} \\
\hline
Power threshold        & received power      & $\mathbb{1}[\,\text{flagged}\,]$              & 3.0  & --   & PPO \\
Kurtosis, Gaussian head& kurtosis            & $\mathbb{1}[\,\text{flagged}\,]$              & 0.5  & --   & PPO \\
Kurtosis, flow head    & kurtosis            & $\max(0,\ \kappa(\mathbf{y}) - \tau)$         & 2.0  & 0.02 & PPO \\
Kurtosis, mixture head & kurtosis            & $\max(0,\ \kappa(\mathbf{y}) - \tau)$         & 2.0  & 0.02 & PPO \\
Kurtosis, direct       & kurtosis            & $\max(0,\ \kappa(\mathbf{y}) - \tau)$         & 2.0  & 0.02 & direct gradient \\
Two agents, direct     & kurtosis            & $\max(0,\ \kappa(\mathbf{y}) - \tau)$         & 2.0  & 0.02 & direct gradient, joint \\
Learned, white-box     & spectrogram CNN     & $-\log(1 - \hat{p} + \epsilon)$               & 0.3  & 0.1  & MAPPO \\
Learned, black-box     & spectrogram CNN     & $-\log(1 - \hat{p} + \epsilon)$               & 0.3  & 0.05 & MAPPO \\
\hline
Minimal model          & Neyman--Pearson     & flagged frames, $\beta$ swept                 & --   & \multicolumn{2}{l}{power is a hard budget} \\
\hline\hline
\end{tabular}
\end{table*}
